This paper has developed a computational framework for distinguishing between decisive shock and strategic exhaustion as mechanisms of war termination. Our two complementary metrics, the Decisive Shock Score (DSS) and the Strategic Exhaustion Score (SES), have been applied to a comprehensive dataset of 4,812 wars, and our dynamical simulation model has been evaluated against 7 historical case studies.

Several findings emerge from this analysis. First, a logistic regression model using ten material-capability features achieves 55.0\% cross-validated accuracy (only 2.7pp above the null baseline), while a random forest achieves 72.7\% accuracy on the same features. This comparison indicates that nonlinear interactions among these specific material-capability features contain more predictive information than linear additive relationships, though both models are limited to the COW National Material Capabilities dataset and should not be interpreted as general tests of war predictability. This demonstrates that warfare dynamics involve complex structural interactions that simple additive models cannot represent. Second, among the 91 wars with complete battle-level data, clear-cut decisive outcomes are rare (2.2\%), with the majority classified as uncertain or negotiated, suggesting that many wars involve mixed or attritional dynamics. Third, our mechanism classifier agrees with author-assigned labels in 6 of 7 cases by separating termination events (how wars end) from dominant mechanisms (why wars become unwinnable). This separation is the paper's core methodological contribution: a war may end because a capital falls, but the reason it became unwinnable may be exhaustion.

These findings suggest that the decisive battle is best understood not as an alternative to attrition but as a mechanism that operates within an attritional context. The question for future research is not whether decisive battles exist (they clearly do) but under what conditions a decisive shock can redirect a war's trajectory versus when it merely is consistent with an outcome that exhaustion has already made inevitable.

Future work should extend the DSS and SES metrics to incorporate non-Western conflicts and pre-modern warfare, develop more sophisticated simulation models that capture alliance dynamics and intelligence operations, and apply our framework to contemporary and near-future conflicts where the balance between precision strikes and prolonged competition is a central strategic question. The ultimate goal is a theory of war termination that integrates both Mahan's decisive battle and Clausewitz's exhaustion into a unified dynamical framework.
